\documentclass[a4paper,10pt]{article}

% Packages
\usepackage[utf8]{inputenc}
\usepackage[T1]{fontenc}
\usepackage{geometry}
\usepackage{amsmath, amssymb, amsfonts, mathtools}
\usepackage{longtable} % For tables spanning multiple pages
\usepackage{booktabs}  % For prettier tables
\usepackage{array}     % For better column definitions
\usepackage{hyperref}

% Page Layout
\geometry{left=2cm, right=2cm, top=2.5cm, bottom=2.5cm}

% Custom Column Types
% L: Left aligned, variable width
\newcolumntype{L}[1]{>{\raggedright\arraybackslash}p{#1}}

\title{\textbf{Mathematical Notation Reference}\\\large Measure Theory, Set Theory, Probability \& Statistics}
\author{Personal Reference}
\date{\today}

\begin{document}

\maketitle
% \tableofcontents
% \newpage

\section{Introduction}
This document serves as a reference for mathematical notations used in Measure Theory, Set Theory, Probability, and Statistics.

% Table Configuration
\setlength{\tabcolsep}{10pt} % Padding between columns
\renewcommand{\arraystretch}{1.5} % Padding between rows

\section{Set Theory}
\begin{longtable}{ | L{0.15\textwidth} | L{0.45\textwidth} | L{0.25\textwidth} | }
    \hline
    \textbf{Notation} & \textbf{Explanation} & \textbf{Purpose} \\
    \hline
    \endfirsthead
    
    \hline
    \textbf{Notation} & \textbf{Explanation} & \textbf{Purpose} \\
    \hline
    \endhead

    \hline
    \endfoot

    \hline
    \endlastfoot

    % Entries start here
    $x \in A$ 
    & Element $x$ belongs to set $A$.
    & Basic set membership. \\
    \hline

    $A \subseteq B$
    & Set $A$ is a subset of set $B$.
    & Indicates inclusion. \\
    \hline

    $A \cap B$
    & Intersection of sets $A$ and $B$ (elements common to both).
    & Logical AND operation on sets. \\
    \hline

    $A \cup B$
    & Union of sets $A$ and $B$ (elements in either $A$ or $B$).
    & Logical OR operation on sets. \\
    \hline

    $\emptyset$
    & The empty set containing no elements.
    & Identity element for union. \\
    \hline

    $2^A$ or $\mathcal{P}(A)$
    & Power set of $A$ (set of all subsets of $A$).
    & Defining sigma-algebras. \\
    \hline
    
    $\forall$ 
    & Universal Quantifier (For all) 
    & Defining properties that must hold for every element (e.g., continuity). \\ \hline

    $\exists$ 
    & Existential Quantifier (There exists) 
    & Asserting the existence of a specific object (e.g., existence of a limit). \\ \hline

    $f^{-1}(B)$ 
    & Inverse Image 
    & The core tool for defining Random Variables (Mapping values back to events). \\ \hline

\end{longtable}


\section{Measure Theory}
\begin{longtable}{ | L{0.15\textwidth} | L{0.45\textwidth} | L{0.25\textwidth} | }
    \hline
    \textbf{Notation} & \textbf{Explanation} & \textbf{Purpose} \\
    \hline
    \endfirsthead

    \hline
    \textbf{Notation} & \textbf{Explanation} & \textbf{Purpose} \\
    \hline
    \endhead

    \hline
    \endfoot

    \hline
    \endlastfoot

    $\mu, \nu$ 
    & General Measures 
    & Generalizing the concept of ``size'' (volume, mass, or count). \\ \hline

    $\lambda$ 
    & Lebesgue Measure 
    & The standard way to measure length/area/volume on the real line $\mathbb{R}$. \\ \hline

    $\mathcal{B}(\mathbb{R})$ 
    & Borel $\sigma$-algebra 
    & The standard collection of ``well-behaved'' sets on $\mathbb{R}$ used in statistics. \\ \hline

    $\nu \ll \mu$ 
    & Absolute Continuity 
    & A condition ensuring that if $\mu$ sees nothing (zero), $\nu$ also sees nothing. \\ \hline

    $\frac{d\mu}{dP}$ 
    & Radon-Nikodym Derivative 
    & Formally defining the Probability Density Function (PDF) as a ratio of measures. \\ \hline

    $\int_A f \, d\mu$ 
    & Lebesgue Integral 
    & Calculating the ``average'' or ``total'' over a set $A$ using measure $\mu$. \\ \hline

\end{longtable}

\section{Probability \& Statistics}
\begin{longtable}{ | L{0.15\textwidth} | L{0.45\textwidth} | L{0.25\textwidth} | }
    \hline
    \textbf{Notation} & \textbf{Explanation} & \textbf{Purpose} \\
    \hline
    \endfirsthead

    \hline
    \textbf{Notation} & \textbf{Explanation} & \textbf{Purpose} \\
    \hline
    \endhead

    \hline
    \endfoot

    \hline
    \endlastfoot

    $\Omega$ 
    & Sample Space 
    & The set of all possible ``raw'' outcomes in an experiment. \\ \hline

    $\mathcal{F}$ 
    & $\sigma$-algebra 
    & The collection of ``measurable'' events or the ``information structure'' available. \\ \hline

    $P$ 
    & Probability Measure 
    & A function that assigns a value in $[0,1]$ to every event in $\mathcal{F}$. \\ \hline

    $\mathbb{1}_A$ 
    & Indicator Function 
    & Converting a set $A$ into a function ($1$ if inside, $0$ if outside) for calculations. \\ \hline

    $X_n \xrightarrow{a.s.} X$ 
    & Almost Sure Convergence 
    & Proving the Strong Law of Large Numbers (the most rigorous convergence). \\ \hline

    $X_n \xrightarrow{P} X$ 
    & Convergence in Probability 
    & The standard convergence for consistent estimators in large samples. \\ \hline

    $X_n \xrightarrow{d} X$ 
    & Convergence in Distribution 
    & The foundation of the Central Limit Theorem (limit behavior of distributions). \\ \hline

    $\theta$ 
    & Parameter 
    & The ``hidden truth'' or the target value in a statistical model. \\ \hline

    $\hat{\theta}_n$ 
    & Estimator 
    & The rule or formula used to guess $\theta$ based on $n$ data points. \\ \hline

    $\mathbb{E}[X|\mathcal{G}]$ 
    & Conditional Expectation 
    & Predicting $X$ given a specific sub-information (sub-$\sigma$-algebra) $\mathcal{G}$. \\ \hline

\end{longtable}

\end{document}
